\documentclass[12pt,a4paper]{article}
\usepackage{expediente}
\begin{document}
% =============================================================================
% FICHA DE DATOS DEL EXPEDIENTE --- actualizado 20 ago 2026
% =============================================================================

% --- VENDEDOR / TITULAR REGISTRAL -------------------------------------------
\newcommand{\VendedorNombres}{RICHAR DARWIN}
\newcommand{\VendedorApellidos}{CUTIRE CONDORI}
\newcommand{\VendedorNombreCompleto}{RICHAR DARWIN CUTIRE CONDORI}
\newcommand{\VendedorDNI}{42610690}
\newcommand{\VendedorEstadoCivil}{soltero, sin unión de hecho vigente}
\newcommand{\VendedorOcupacion}{\faltante{[ocupación de Richar]}}
\newcommand{\VendedorDomicilio}{\faltante{[domicilio actual de Richar]}}
\newcommand{\VendedorTelefono}{\faltante{[teléfono de Richar]}}
\newcommand{\VendedorCorreo}{\faltante{[correo]}}
\newcommand{\VendedorNacionalidad}{peruano}
\newcommand{\VendedorFechaNac}{\faltante{[fecha de nacimiento]}}
\newcommand{\VendedorDNIEmision}{\faltante{[fecha de emisión del DNI de Richar --- la pide SUNARP]}}

% --- CONVIVIENTE: NO CORRESPONDE -------------------------------------------
\newcommand{\ConvivienteNombre}{NO CORRESPONDE}
\newcommand{\ConvivienteDNI}{NO CORRESPONDE}
\newcommand{\ConvivienteDesde}{NO CORRESPONDE}
\newcommand{\ConvivienteInscrita}{no existe unión de hecho vigente}
\newcommand{\ConvivienteDomicilio}{NO CORRESPONDE}

% --- COMPRADORES (copropiedad en partes iguales) ---------------------------
\newcommand{\CompradorUnoNombre}{NICO ALVARO CUTIRE ARCE}
\newcommand{\CompradorUnoDNI}{48325017}
\newcommand{\CompradorUnoEmision}{01 de enero de 2024}
\newcommand{\CompradorUnoEstadoCivil}{soltero}
\newcommand{\CompradorUnoOcupacion}{ingeniero informático}
\newcommand{\CompradorUnoDomicilio}{Jr.\ Wiracocha, Sicuani, provincia de Canchis, Cusco, Perú}
\newcommand{\CompradorUnoTelefono}{984715108}

\newcommand{\CompradorDosNombre}{LUCY CUTIRE ARCE}
\newcommand{\CompradorDosDNI}{47478852}
\newcommand{\CompradorDosEmision}{09 de enero de 2026}
\newcommand{\CompradorDosEstadoCivil}{soltera}
\newcommand{\CompradorDosOcupacion}{contador}
\newcommand{\CompradorDosDomicilio}{Comunidad Hanocca, distrito de Layo, provincia de Canas, Cusco, Perú}
\newcommand{\CompradorDosTelefono}{953440875}

\newcommand{\CompradorNombres}{NICO ALVARO y LUCY}
\newcommand{\CompradorApellidos}{CUTIRE ARCE}
\newcommand{\CompradorNombreCompleto}{NICO ALVARO CUTIRE ARCE y LUCY CUTIRE ARCE}
\newcommand{\CompradorDNI}{48325017 y 47478852}
\newcommand{\CompradorEstadoCivil}{solteros}
\newcommand{\CompradorConyuge}{NO CORRESPONDE}
\newcommand{\CompradorConyugeDNI}{NO CORRESPONDE}
\newcommand{\CompradorOcupacion}{ingeniero informático y contador, respectivamente}
\newcommand{\CompradorDomicilio}{Sicuani (Canchis) y Layo (Canas), Cusco}
\newcommand{\CompradorTelefono}{984715108 / 953440875}
\newcommand{\CompradorCorreo}{\faltante{[correo de contacto]}}
\newcommand{\CompradorNacionalidad}{peruanos}
\newcommand{\PorcentajeCopropiedad}{50\,\% cada uno}

% --- INTERMEDIARIO / OCUPANTE DEL REMANENTE --------------------------------
\newcommand{\IntermediarioNombre}{LUCIO GIL CUTIRE PARI}
\newcommand{\IntermediarioDNI}{42467850}
\newcommand{\IntermediarioEstadoCivil}{\faltante{[estado civil de Lucio]}}
\newcommand{\IntermediarioOcupacion}{\faltante{[ocupación]}}
\newcommand{\IntermediarioDomicilio}{\faltante{[domicilio de Lucio]}}
\newcommand{\IntermediarioTelefono}{\faltante{[teléfono de Lucio]}}
\newcommand{\IntermediarioParentesco}{familiar (apellido Cutire)}

% --- PREDIO MATRIZ ----------------------------------------------------------
\newcommand{\PartidaMatriz}{\faltante{[N.º de partida electrónica SUNARP --- pendiente]}}
\newcommand{\OficinaRegistral}{Oficina Registral de Puerto Maldonado}
\newcommand{\AreaMatriz}{\faltante{[área total inscrita; las partes refieren del orden de 100 ha]}}
\newcommand{\UnidadCatastral}{\faltante{[código catastral / UC, si existe]}}
\newcommand{\NombreFundo}{\faltante{[nombre del fundo, si tiene]}}
\newcommand{\KmCarretera}{Interoceánica Sur, al este del Puente Jayave (el puente queda al oeste del predio, sobre el río; no debajo del lote)}

% --- PREDIO A INDEPENDIZAR --------------------------------------------------
\newcommand{\AreaIndependizarHa}{10{,}00}
\newcommand{\AreaIndependizarMetros}{100 000{,}00}
\newcommand{\PerimetroM}{2 200{,}58}
\newcommand{\FrenteM}{100{,}29}
\newcommand{\FondoM}{1 000{,}00}

\newcommand{\PUnoLat}{$-12{,}912452$}
\newcommand{\PUnoLon}{$-70{,}170058$}
\newcommand{\PDosLat}{$-12{,}912497$}
\newcommand{\PDosLon}{$-70{,}170981$}
\newcommand{\PTresLat}{$-12{,}903469$}
\newcommand{\PTresLon}{$-70{,}171438$}
\newcommand{\PCuatroLat}{$-12{,}903424$}
\newcommand{\PCuatroLon}{$-70{,}170515$}

\newcommand{\PUnoE}{373 065{,}00}
\newcommand{\PUnoN}{8 572 256{,}00}
\newcommand{\PDosE}{372 965{,}00}
\newcommand{\PDosN}{8 572 251{,}00}
\newcommand{\PTresE}{372 911{,}00}
\newcommand{\PTresN}{8 573 249{,}00}
\newcommand{\PCuatroE}{373 011{,}00}
\newcommand{\PCuatroN}{8 573 255{,}00}

\newcommand{\FechaPago}{15 de agosto de 2018}
\newcommand{\MontoPago}{S/\ 20 000{,}00}
\newcommand{\MontoPagoLetras}{VEINTE MIL Y 00/100 SOLES}
\newcommand{\FormaPagoHistorica}{dinero en efectivo, sin medio de pago del sistema financiero}

% Nombres genéricos --- sustituir por los reales al ir a notaría
\newcommand{\NotarioNombre}{Dr.\ JUAN CARLOS RÍOS SALINAS (nombre genérico a sustituir)}
\newcommand{\NotariaCiudad}{Puerto Maldonado}
\newcommand{\ColegiaturaAbogado}{CAL Madre de Dios N.º 0000 (genérico)}
\newcommand{\AbogadoNombre}{Abog.\ ANA MARÍA QUISPE HUAMÁN (nombre genérico a sustituir)}
\newcommand{\IngenieroNombre}{Ing.\ PEDRO ANTONIO MAMANI CONDORI (nombre genérico a sustituir)}
\newcommand{\IngenieroCIP}{CIP 000000 (genérico)}
\newcommand{\Municipalidad}{Municipalidad Distrital de Inambari}
\newcommand{\Provincia}{Tambopata}
\newcommand{\Departamento}{Madre de Dios}
\newcommand{\Distrito}{Inambari}

\newcommand{\LugarOtorgamiento}{Puerto Maldonado}
\newcommand{\DiaOtorgamiento}{\faltante{[día]}}
\newcommand{\MesOtorgamiento}{\faltante{[mes]}}
\newcommand{\AnioOtorgamiento}{2026}

\newcommand{\TestigoUno}{\faltante{[testigo 1: nombres, DNI, domicilio]}}
\newcommand{\TestigoDos}{\faltante{[testigo 2: nombres, DNI, domicilio]}}

\InicioDocumento{Hoja de vértices y coordenadas de trabajo --- polígono de 10,00 ha}{DOC-08}

\noindent Este documento fija los cuatro vértices usados para visualizar el recorte. \textbf{Datum de trabajo:} WGS84. \textbf{Proyección de apoyo:} UTM zona 19 Sur (EPSG:32719), calculada geodésicamente a partir de los dos puntos Google aportados por las partes. \textbf{No es plano de independización.} El profesional \IngenieroNombre, CIP \IngenieroCIP, sustituirá estas cifras por las de campo.

\Clausula{Criterio geométrico}
Se tomó como frente el segmento de cien metros entre P1 y P2, coincidente con la calzada de la \textbf{Interoceánica Sur}. Yendo Mazuco $\rightarrow$ Puerto Maldonado (rumbo aproximado este, azimut $\approx 87^\circ$), la \textbf{izquierda} es el norte. Se proyectó un rectángulo de $\FrenteM$\,m de frente por $\FondoM$\,m de fondo, equivalente a $\AreaIndependizarHa$ hectáreas.

\Clausula{Cuadro de vértices}
{\small\renewcommand{\arraystretch}{1.32}
\noindent\begin{tabularx}{\textwidth}{>{\bfseries}c >{\raggedright\arraybackslash}p{3.3cm} c c c c}
\toprule
Vértice & Rol & Latitud & Longitud & Este UTM 19S & Norte UTM 19S \\
\midrule
P1 & Frente este (dato) & \PUnoLat & \PUnoLon & \PUnoE & \PUnoN \\
P2 & Frente oeste (dato) & \PDosLat & \PDosLon & \PDosE & \PDosN \\
P3 & Fondo oeste (calc.) & \PTresLat & \PTresLon & \PTresE & \PTresN \\
P4 & Fondo este (calc.) & \PCuatroLat & \PCuatroLon & \PCuatroE & \PCuatroN \\
\bottomrule
\end{tabularx}}

\vspace{0.45em}
\Clausula{Medidas resultantes}
\renewcommand{\arraystretch}{1.25}
\noindent\begin{tabularx}{\textwidth}{>{\bfseries}p{6.2cm}X}
P1--P2 (frente / sur) & \FrenteM\ m \\
P2--P3 (oeste) & \FondoM\ m \\
P3--P4 (fondo / norte) & \FrenteM\ m \\
P4--P1 (este) & \FondoM\ m \\
Perímetro & \PerimetroM\ m \\
Área & \AreaIndependizarMetros\ m$^2$ \quad $=$ \quad \AreaIndependizarHa\ ha \\
Sentido del polígono & P1 $\rightarrow$ P2 $\rightarrow$ P3 $\rightarrow$ P4 $\rightarrow$ P1 \\
\end{tabularx}

\Clausula{Linderos de trabajo (colindancias a confirmar en campo)}
\begin{itemize}[leftmargin=1.15cm]
\item \textbf{Por el Sur:} vía nacional Interoceánica Sur (PE-30C), frente \FrenteM\ m, entre P1 y P2. El Puente Jayave queda \textbf{al oeste de P2}, sobre el río, no bajo el predio. Faja de derecho de vía a cargo del MTC.
\item \textbf{Por el Oeste:} línea P2--P3 de \FondoM\ m, colindante con el \textbf{remanente} del predio matriz de \VendedorNombreCompleto, u otro que el plano de verificador determine.
\item \textbf{Por el Norte:} línea P3--P4 de \FrenteM\ m, colindante con el remanente del predio matriz.
\item \textbf{Por el Este:} línea P4--P1 de \FondoM\ m, colindante con el remanente del predio matriz.
\end{itemize}

\Clausula{Observaciones técnicas de las partes}
\begin{enumerate}[leftmargin=1.15cm]
\item Si el fundo útil está al \textbf{sur} de la carretera y no al norte, este cuadro queda sin efecto y debe recalcularse.
\item Google Maps tiene error típico de varios metros. Los hitos de campo priman sobre estas coordenadas.
\item El profesional deberá amojonar P1 a P4 (o los vértices definitivos) con hitos de concreto u otro material duradero, y consignar el factor de escala y el origen de la red geodésica.
\end{enumerate}

\Clausula{Declaración}
Las partes reconocen estos vértices \textbf{solo como hipótesis de trabajo} para minutas y plano esquemático, y se obligan a aceptar el área, linderos y coordenadas que resulten del levantamiento formal, siempre que se mantenga el objeto de independizar aproximadamente diez hectáreas en el frente indicado.

\LugarFecha
\TresFirmas{EL TITULAR}{\VendedorNombreCompleto\\DNI \VendedorDNI}{EL COMPRADOR}{\CompradorNombreCompleto\\DNI \CompradorDNI}{EL INTERMEDIARIO}{\IntermediarioNombre\\DNI \IntermediarioDNI}

\vspace{0.7em}
\noindent
\BloqueFirmaSimple{PROFESIONAL QUE SUSTITUIRÁ ESTE CUADRO}{\IngenieroNombre\\CIP \IngenieroCIP}

\end{document}
